\begin{figure*}[t]
    \centering
      \includegraphics[width=1\textwidth]{model2.pdf}
    \caption{Given an image and human instruction as the input, we introduce \textit{GPT4-Assisted Visual Instruction Evaluation (GAVIE)} to assess the output from current LMMs, such as \textit{MiniGPT4 and LLaVA}. \textbf{\textcolor{blue}{BLUE}} represents LMMs can not accurately follow human instructions while \textbf{\textcolor{red}{RED}} means they suffer from the hallucination problem. We can generate a more robust answer after finetuning MiniGPT4 on our proposed \textit{LRV-Instruction} dataset.  }
    \label{fig:model}
    \vspace{-0.15in}
\end{figure*}
% The recent advancements in natural language processing literature have seen a significant development of models capable of comprehending and following instructions given in natural language.
Significant progress has been made in the field of natural language processing, leading to the development of models that can comprehend and follow instructions given natural language inputs~\cite{wang2022self, gong2023multimodal, ouyang2022training, brown2020language}. These models harness the power of large language models (LLM) and rely on high-quality instruction data. 
% Similar attempts have been introduced in multi-modal models as well. 
Similarly, efforts have been made to introduce similar capabilities to multi-modal models.
GPT4 \cite{peng2023instruction} has demonstrated impressive performance in multi-modal conversations with humans, yet the techniques contributing to its extraordinary capabilities remain opaque. As a result, several large multi-modal models (LMMs) have recently emerged \cite{zhu2023minigpt, liu2023visual, gong2023multimodal, dai2023instructblip}, such as MiniGPT4 \cite{zhu2023minigpt} and LLaVA \cite{liu2023visual}, both utilize the Vicuna \cite{chiang2023vicuna} as the language generator but with different vision encoders \cite{radford2021learning, li2022blip}. InstructBLIP \cite{dai2023instructblip} is initialized from a pre-trained BLIP-2 \cite{li2023blip} while Multimodal-GPT (MMGPT) \cite{gong2023multimodal} is built on Flamingo \cite{alayrac2022flamingo, anas_awadalla_2023_7733589}.

A recent study~\cite{li2023otter}   revealed that the hallucination issue of LLM, although not desired, is inherited by these LMMs \cite{zhu2023minigpt, liu2023visual, gong2023multimodal, dai2023instructblip}. %Hallucination \cite{bang2023multitask} is one of the major ethical concerns of LLMs, and it can have harmful consequences  
Hallucination, a major ethical concern associated with LLMs~\cite{bang2023multitask}, can lead to harmful consequences,
especially when users without adequate domain knowledge over-rely on these increasingly convincing language models.
In the context of LMM hallucinations, the model can generate descriptions with conflicting information to the given image. For instance, as shown in Fig. \ref{fig:model}~(highlighted in red), existing LMMs \cite{zhu2023minigpt, liu2023visual, dai2023instructblip, gong2023multimodal} tend to describe nonexistent objects such as a "dog" engaging in a nonexisting activity like "playing with a ball".
Additionally, the model may generate long image descriptions without following human instructions (highlighted in blue). 

\textbf{\textit{What are the likely causes of these hallucinations?}} As current LMMs are built on strong LLMs, they may over-rely on language priors and generate words more likely to go together with the instruction text regardless of the image content. What's more, LMMs, such as MiniGPT4 \cite{zhu2023minigpt} and LLaVA \cite{liu2023visual}, employ synthetic instruction data for training,  which are generally long and involve nonexistent objects, activities, or relationships in the image.

\textbf{\textit{Why can't LMMs accurately follow human instructions?}} We conjecture it is due to the lack of diversity in their training data. For example, MiniGPT4 \cite{zhu2023minigpt} is only instruct tuning with four instruction templates designed for image captioning tasks. Though MMGPT \cite{gong2023multimodal} and InstructBLIP \cite{dai2023instructblip} combine several datasets as the instruction tuning data, their instructions and answers are still based on a few templates.

To address these challenges, we present \textit{LRV-Instruction}, a large and diverse visual instruction benchmark. Our benchmark consists of 120k visual instructions generated by GPT4, taking inspiration from the success of recent GPT models in text-annotation tasks~\cite{liu2023gpteval}. Unlike previous studies that focused on limited tasks and pre-defined templates created by human experts~\cite{zhu2023minigpt, dai2023instructblip, gong2023multimodal}, \textit{LRV-Instruction} covers 16 vision-language tasks with open-ended instructions and answers, as shown in Fig.~\ref{fig:tasks} and Tab.~\ref{tab:stat1}. As observed by \cite{li2023evaluating}, current LMMs tend to answer \textit{"Yes"} for any instructions presented to the model, even when the proper answer should be \textit{"No"}. Some models even answer with \textit{"Yes"} for over 90\% of the negative testing examples during their evaluation. Our investigation reveals that most LMMs are finetuned on unbalanced datasets containing only positive instructions (Tab. \ref{tab:comparison}). To enable LMMs to respond to human instructions more faithfully, we design \textit{LRV-Instruction} to include both negative and positive instructions for robust instruction tuning. Our negative instructions are generated at two semantic levels (Fig. \ref{fig:neg_example}): \textit{(i) Nonexistent Element Manipulation} and \textit{(ii) Existent Element Manipulation} in two different formats, \textit{Declarative} and \textit{Interrogative}. To improve the robustness and flexibility of the evaluation on visual instruction tuning, we propose \textit{GPT4-Assisted Visual Instruction Evaluation (GAVIE)} to assess the LMM output in two different aspects: \textit{Relevancy} to evaluate the instruction-following performance and \textit{Accuracy} to measure the visual hallucination in the LMM output. \textit{GAVIE} does not require human-annotated groundtruth answers \cite{rohrbach2018object} and can be easily adapted to different instruction formats instead of specific designs in \cite{li2023evaluating}.

% Based on our experiments,
We empirically evaluate four publicly available LMMs~\cite{zhu2023minigpt,liu2023visual,dai2023instructblip,gong2023multimodal} on our benchmark and found that existing LMMs seriously hallucinate when prompted with our negative instructions. We further verify the effectiveness of our \textit{LRV-Instruction} by finetuning MiniGPT4~\cite{zhu2023minigpt} on this more balanced data. Our instruct-tuned MiniGPT4 suffers much less from hallucination and achieves state-of-the-art performance compared to the original MiniGPT4, LLaVA \cite{liu2023visual}, InstructBLP \cite{dai2023instructblip} and MMGPT \cite{gong2023multimodal} on both our evaluation set and public benchmarks \cite{li2023evaluating, hudson2019gqa}. We also observe that \textit{Existent Element Manipulation} instructions are more challenging than \textit{Nonexistent Element Manipulation} instructions for LMMs. Furthermore, a robust model performance requires a balanced ratio between positive and negative instances. To sum up, our contributions are three-fold:
\begin{itemize}[nosep,leftmargin=*]
    \item We build \textit{LRV-Instruction}, a large and diverse dataset containing 120k visual instructions, with 16 vision and language tasks and negative instructions in different semantic levels and styles.
    \item We propose \textit{GAVIE}, a novel approach to evaluate visual instruction tuning without requiring groundtruth answers and pre-designed instruction formats.
    \item  We conduct comprehensive experiments to investigate the hallucination of current LMMs. The empirical study validates the effectiveness of \textit{LRV-Instruction} for robust visual instruction tuning.
\end{itemize}

\begin{figure*}[t]
    \centering
      \includegraphics[width=1\textwidth]{task2.pdf}
    \caption{Examples of negative instructions in our \textit{LRV-Instruction} dataset and output from current LMMs and our model. \textbf{\textcolor{red}{RED}} means inconsistent elements in the negative instructions. }
    \label{fig:neg_example}
    \vspace{-0.15in}
\end{figure*}



